\chapter{Conclusion and Future Scope}
The fundamental displacement of synchronous, high-inertia machines by Inverter-Based Resources (IBRs) fundamentally breaks classical power system protection paradigms. Because IBRs dictate fault current amplitude and phase through power electronic converter controls and Low Voltage Ride Through (LVRT) mandates, they exhibit non-linear, current-limited fault signatures that systematically blind traditional overcurrent, distance, and sequence-polarized directional relays. 
\section{Core Technical Contributions}
This thesis aggressively confronts these deterministic vulnerabilities, engineering advanced mathematical and data-driven methodologies to reclaim protection dependability, selectivity, and speed in highly dynamic, converter-dominated networks.
\subsection{The Enhanced Voltage Protection Module (EVPM)} 
The primary deterministic breakthrough of this research is the formulation of the EVPM. Bypassing the crippling limitations of conventional RMS-based undervoltage algorithms, the EVPM executes a geometric time-domain characterization of voltage sags utilizing the Space Phasor Model (SPM). By applying direct least squares estimation to fit the dynamic, three-phase voltage locus to an ellipse, the algorithm precisely extracts the Characteristic Voltage (CV), Positive Negative Factor (PNF), and tilt angle (ϕ). This physics-based parameterization empowers the relay to instantaneously differentiate between genuine fault conditions and non-fault transient sags, achieving highly selective fault-phase identification and seamless coordination with upstream reclosers without violating mandatory LVRT grid codes.
\subsection{Four-Level Microgrid Protection Architecture}
Addressing the distinct topological challenges of microgrids—specifically the severe variance in short-circuit capacities and bidirectional flows between grid-connected and islanded modes—this thesis advances a robust Four-Level Protection Architecture. This hierarchical strategy successfully synthesizes localized voltage-restrained overcurrent elements (Level 0), ultra-low latency peer-to-peer IEC 61850 GOOSE-enabled differential interlocking (Level 1), and centralized supervisory adaptive setting controls (Level 2). By accommodating variable fault levels and zero-sequence isolation constraints, this layered framework guarantees absolute microgrid selectivity regardless of the operational state.
\subsection{Transform-Domain AI/ML Integration}
To resolve the vulnerabilities of curve-fitting estimations in high-noise, transient-heavy environments, this work introduces an advanced AI-driven augmentation: the Transform-Domain Deep Convolutional Neural Network (CNN). By processing the 2-D and 3-D hierarchical representations of the SPM under fault conditions, the CNN-based architecture classifies complex fault events independent of sampling frequency or duration. This delivers a highly accurate fault classification matrix that remains impervious to the highly distorted, reactive-current waveforms injected by modern Type III and Type IV wind and solar converters.
\section{Future Scope and Trajectory}
The future of active distribution network protection demands an escalation from isolated algorithmic solutions to ubiquitous, system-wide dynamic protection networks. Future engineering efforts must pivot toward:
\subsection{Physics-Informed Neural Networks (PINNs)}
Evolving the current deep learning models into PINNs to create self-learning, adaptive relays capable of real-time boundary adjustment, seamlessly integrating power system dynamics with data-driven inference to distinguish between complex switching transients and high-impedance faults.
\subsection{Active Smart Inverter Coordination}
Leveraging the inherent Volt-Var and Volt-Watt control functionalities of distributed smart inverters, optimized via Non-Sorted Genetic Algorithms (NSGA), to actively suppress voltage violations and manage fault currents directly at the Distribution Management System (DMS) layer.
\subsection{Digital-Twin Validation and Cyber-Resilience}
Expanding real-time Hardware-in-the-Loop (HIL) testing platforms to aggressively stress-test these AI/ML models. This includes validating algorithmic resilience against packet latency, communication failures in IEC 61850 networks, and hardening the protection architectures against sophisticated cyber-attacks or false data injections.